\section{Giới thiệu}
\label{sec:introduction}

Ứng dụng ngân hàng số cung cấp các chức năng truy cập tài khoản, giao dịch, xác
thực và quản trị thông qua các hệ thống hướng web. Những hệ thống này phải trả
lời hai câu hỏi bổ sung cho nhau: người dùng có được phép thực hiện một hành
động hay không, và chuỗi sự kiện quan sát được có biểu hiện hành vi đáng ngờ hay
không. RBAC giải quyết câu hỏi thứ nhất bằng cách ánh xạ người dùng với vai trò
và vai trò với quyền. Phân tích nhật ký bảo mật giải quyết câu hỏi thứ hai bằng
cách lưu giữ bằng chứng sự kiện phục vụ phát hiện, điều tra và phản ứng. Các
hướng dẫn về quản lý log cũng coi việc tổ chức thu thập và phân tích nhật ký là
nền tảng cho điều tra sự cố \citep{Kent2006SP80092}.

Động lực của nghiên cứu này xuất phát từ thực tế rằng trong các nguyên mẫu nhỏ,
hai nguồn thông tin trên thường được xử lý tách biệt: logic kiểm soát truy cập
chặn hoặc cho phép yêu cầu, còn phân tích nhật ký phát hiện mẫu mà không xét đến
ngữ cảnh vai trò. Sự tách rời này làm giảm khả năng giải thích của cảnh báo. Ví
dụ, một thao tác bị từ chối có thể là bình thường nếu người dùng không có quyền,
nhưng nhiều lần bị từ chối từ một IP mới sau một số lần đăng nhập thất bại có
thể cho thấy hành vi thăm dò. Ngược lại, một hành động được RBAC cho phép vẫn có
thể rủi ro nếu diễn ra ngoài ngữ cảnh hành vi thông thường.

Bài báo này trình bày thiết kế và nguyên mẫu RBAC Guard cho bài toán trên.
Trọng tâm là mô hình RBAC xác định và pipeline phân tích log theo luật: parser
chuẩn hóa event, bộ máy luật đánh giá chỉ báo SQL injection và dò đoán mật khẩu,
còn CSDL RBAC cung cấp cơ sở để đối chiếu quyền khi log ghi nhận truy cập hoặc
ủy quyền. Thiết kế tạo alert giữ liên kết với event nguồn, luật và evidence.
Các tín hiệu hành vi theo ngữ cảnh và gom nhóm sự cố là phần mở rộng tùy chọn để
triage; chúng không phải cơ chế cấp quyền thích nghi hay mô hình học máy. Cách
tiếp cận này phù hợp với việc xem RBAC là một nguồn ngữ cảnh có thể kiểm tra
trong giám sát truy cập \citep{Rose2020SP800207,Xiao2022ContextRisk}.

\subsection{Đóng góp}
\label{subsec:contributions}

Bài báo có các đóng góp sau:

\begin{itemize}
  \item Một mô hình tích hợp RBAC--phân tích log, xác định ranh giới giữa quyết
  định phân quyền, event đầu vào không tin cậy và bằng chứng alert đầu ra.
  \item Hai luật phát hiện xác định, có thể giải thích cho chỉ báo SQL injection
  và chuỗi dò đoán mật khẩu; kiểm tra quyền RBAC được dùng để bổ sung bằng chứng
  cho event truy cập/ủy quyền.
  \item Một nguyên mẫu có thể tái lập trên dataset tổng hợp, xuất alert và
  metadata để kiểm chứng đường truy vết từ kết quả về luật và event nguồn.
\end{itemize}

\subsection{Câu hỏi nghiên cứu}
\label{subsec:research-questions}

Nghiên cứu được tổ chức quanh các câu hỏi sau:

\begin{description}
  \item[RQ1] Làm thế nào để thiết kế mô hình RBAC kết hợp với pipeline phân tích
  nhật ký có khả năng truy vết trong ngữ cảnh ứng dụng ngân hàng?
  \item[RQ2] Làm thế nào để biểu diễn rủi ro SQL injection và dò đoán mật khẩu
  thành các luật phát hiện xác định, có evidence kiểm tra được?
  \item[RQ3] Nguyên mẫu có tái lập được các kịch bản tổng hợp đã thiết kế và liên
  kết alert với event, luật và trạng thái RBAC liên quan hay không?
\end{description}
